\chapter{Results, Evaluation, and Discussion}
\label{chapter4}

% ──────────────────────────────────────────────────────────────────────────────────────────────────────────
\section{Performance Results}

Performance was measured using the built-in performance logger, which records per-frame time, FPS, mesh
triangle count, and scene rebuild duration to a CSV file for analysis. The two subsections below cover the
primary variables under direct implementation control: LOD tier and geometry allocation strategy.

\subsection{LOD Effectiveness}

Table~\ref{tab:lod-perf} shows the effect of LOD tier on mesh index count, average frame time over 30 seconds
and rebuild time for the meadow scene at year 3 with the seed set to 13. The three mesh-producing tiers
(Near, Mid, Far) are defined in Table~\ref{tab:lod}; LOD boundaries were altered so that all plants fell into
one tier (confirmed through the debug UI and performance logger output), frustum culling and vsync were
disabled and the camera was set to orbit mode with auto orbit enabled, all other settings were kept as their
defaults.

\begin{table}[h]
  \centering
  \begin{tabular}{|l|r|r|r|}
    \hline
    \textbf{LOD tier} & \textbf{Mesh triangles} & \textbf{Average frame time (ms)} & \textbf{Rebuild time (ms)} \\
    \hline
    Near              & 5816578                 & 5.4961                           & 275.44 \\
    Mid               & 3690514                 & 4.0616                           & 163.82 \\
    Far               & 2197802                 & 2.8490                           & 133.22 \\
    \hline
  \end{tabular}
  \caption{Impact of LOD tier on performance for the meadow scene.}
  \label{tab:lod-perf}
\end{table}

Across the three tiers, mesh triangle count falls by a factor of 2.65 (Near to Far). GPU frame time falls by
a factor of 1.93 over the same range, as the draw call overhead and vertex transform cost have a fixed
component independent of triangle count. All three tiers remain well below the 16.7\,ms threshold for 60\,FPS
interactive use, achieving 182\,FPS at Near and 351\,FPS at Far.

\medskip
Rebuild time falls by a factor of 2.07 from Near to Far, tracking the triangle-count reduction more closely
than frame time does as rebuild time is dominated by CPU geometry generation, which scales with the amount of
geometry produced. The rebuild costs represent a one-time per-change expense caused by switching LOD tiers.
Rebuilds are only triggered by parameter changes, season threshold crossings, or LOD tier transitions,
meaning that the cost is not paid every frame.

\subsection{Allocation and Hash Table Optimisations}

Three allocation-level optimisations were applied to the turtle geometry path to reduce per-plant overhead
during scene rebuilds.

\medskip
First, the voxel occupancy sets use an inline FNV-1a hash function (implemented in approximately ten lines of
Rust, with no external dependency) in place of Rust's default SipHash hasher. SipHash is designed to resist
hash-flooding attacks, a property that is unnecessary for integer-keyed internal data
structures~\cite{aumassonSipHashFastShortInput2012}. For \texttt{(i32,i32,i32)} keys, FNV-1a reduces lookup
and insert cost by a factor of approximately three to five compared to the default hasher, because the hash
computation over twelve bytes of already-uniform integer data requires only twelve XOR-multiply iterations
rather than SipHash's more expensive key schedule~\cite{fowlerFNVNonCryptographicHash2026}.

\medskip
Second, the turtle's internal per-plant buffers (push stack, path vertices, colours, and path entries) are
cleared and reused across plants rather than reallocated. The \texttt{reset()} method calls \texttt{clear()}
followed by \texttt{push()} instead of constructing new \texttt{Vec} instances, preserving the heap capacity
allocated during the previous plant's build for reuse by the next.

\medskip
Third, mesh geometry is moved out of the turtle into the plant cache via \texttt{std::mem::take} rather than
cloned. This eliminates one full copy of six geometry vectors (vertices, normals, colours, indices, UVs,
tangents) per plant per scene rebuild, reducing peak memory allocation during rebuilds proportionally to scene size.

% ──────────────────────────────────────────────────────────────────────────────────────────────────────────
\section{Evaluation Against Objectives}

\subsection{Objective 1: L-System Interpreter}

The interpreter correctly implements all seven rule variants. This is confirmed by the unit test suite
(Section~\ref{sec:testing}), which covers the D0L algae sequence, terminal pass-through, stochastic
probability boundaries, and all forms of context-sensitive matching including branch-skipping and
multi-branch right-context search. Context-sensitive priority was further validated manually for the
\textit{Mycelis} grammar (Figure~\ref{fig:mycelis}), where the acropetal signal chain depends on context
rules firing before general rules for the same symbol.

\paragraph{Limitations:}
Grammars are defined entirely in Rust code: adding a new species requires writing a new \texttt{.rs} file,
implementing the \texttt{Plant} trait, and recompiling. There is no runtime grammar definition language,
which limits the system's use as an authoring tool for non-programmers.

\subsection{Objective 2: Plant Grammars}

Eight distinct plant grammars are implemented across four grammar classes. Plants based on grammars described
in \textit{The Algorithmic Beauty of Plants} (\textit{Lychnis}, \textit{Mycelis}, and \textit{Capsella}) are
visually similar to their respective figures. Other plants produce convincing plant forms with appropriate
branching structure and leaf distribution.

\paragraph{Limitations.}
Quantitative comparison of branch angles with measured plant data has not been performed. Whilst possible,
L-Systems are not used to generate leaves or flowers, a leaf texture is used instead.

\subsection{Objective 3: Environmental Response}

Phototropism, gravitropism, wind, and season response are all implemented and interactive. The season system
demonstrates visually distinct states at spring, summer, autumn, and winter on two levels: a global colour
tint applied per-fragment without a geometry rebuild, and per-species dormancy behaviour
(Figure~\ref{fig:seasons}). Deciduous species (bush and tree) drop their leaves while retaining branch structure.
Herbaceous species (wildflower, fern, mint, \textit{lychnis}, \textit{mycelis}) fade to dormancy by scaling
the effective L-System age. The tree progressively increases leaf-cluster shedding probability. Tropism
produces convincing branch bending toward the light source at moderate strength values.

\paragraph{Limitations.}
The environmental response is applied as a post-processing deflection of the turtle frame rather than as
feedback into the L-System rewriting process. This means the grammar cannot respond to local environmental
conditions (e.g.\ a bud entering shadow and aborting). A full open L-System
implementation~\cite{mechVisualModelsPlants1996} would require the environment to be sampled at each symbol
position during rewriting, significantly increasing implementation complexity. Phototropism is applied
uniformly to all branches of all plants in a scene, it does not model per-branch light levels. In a real
forest, the canopy shades lower branches, and the tree architecture adapts accordingly. Implementing
per-branch irradiance would require either a shadow map or a simplified canopy model.

\subsection{Objective 4: Interactive 3D Visualisation}

The system maintains interactive frame rates across all provided scenes. The egui GUI provides real-time
control of all parameters with immediate visual feedback (Appendix~\ref{app:gui}). The debug HUD displays
FPS, frame time, and total index count with a rolling 128-frame bar graph. The LOD system adapts
automatically to camera position, and the season auto-advance feature allows unattended demonstration of the
full growth cycle.

\paragraph{Limitations.}
The renderer uses a single Lambert diffuse lighting model with no shadow casting or receiving. Plants do
not shade each other, meaning large scenes appear uniformly lit regardless of canopy density. Implementing
even an approximate ambient occlusion term would significantly improve visual realism at the cost of
additional render passes. There is no anti-aliasing. At typical viewing distances branch silhouettes exhibit
temporal aliasing as the camera moves. Adding MSAA would require configuring a multisampled render target and
resolve pass, which was deferred as a lower-priority rendering quality concern.

% ──────────────────────────────────────────────────────────────────────────────────────────────────────────
\section{Conclusions and Future Work}

\subsection{Summary of Contributions}

This project has delivered a real-time procedural vegetation system in Rust with the following novel
contributions relative to existing tools:

\begin{itemize}
  \item A single unified L-System interpreter supporting four rule classes (deterministic, stochastic,
    parametric, and context-sensitive) in a type-safe, zero-overhead Rust implementation.
  \item Eight botanically motivated plant grammars spanning four grammar classes, with per-branch stochastic
    variation and age-based colour.
  \item An integrated environmental response system combining phototropism, gravitropism, wind and a seasonal
    cycle with adaptive rebuild scheduling. Per-species dormancy, leaf-drop, and flowering phenology are
    encoded directly in each grammar's parametric closures and activated by the season value passed through
    \texttt{PlantEnvironment} to \texttt{Plant::actions()}.
  \item A kernel-based ecosystem placement system supporting multi-species communities with weighted
    abundance, self-thinning, and ecological succession, with per-species ecological properties encoded at
    the type level.
  \item A five-tier LOD system (\texttt{Near}/\texttt{Mid}/\texttt{Far}/\texttt{Beyond}/\texttt{Culled}) with
    per-tier leaf-skip probabilities, view-frustum culling via the Gribb--Hartmann method, and an adaptive
    geometry-budget pressure mechanism that dynamically tightens LOD thresholds when assembled geometry
    exceeds the GPU buffer limit.
\end{itemize}

\subsection{Future Work}

\paragraph{Open L-Systems.}
The most significant technical extension would be implementing full open L-System
semantics~\cite{mechVisualModelsPlants1996}, in which the L-System queries the environment (light levels,
occupied space) during rewriting. This would enable branch abortion in shade, root systems competing for
nutrients, and graft compatibility modelling.

\paragraph{Grammar serialisation.}
Currently, all grammars are hard-coded in Rust. A grammar definition language (similar to
L-Py~\cite{boudonLPyLSystemSimulation2012}) would allow users to define and share new plant types without recompiling.

\paragraph{Procedural bark and branch textures.}
Branch cylinders are currently shaded with a flat vertex colour interpolated by age. Procedural or
texture-mapped bark would substantially improve close-up visual quality without changing the grammar or
geometry pipeline.
